\chapter{System Implementation}
\label{chap:implementation}

\section{Implementation Overview}

The proposed methodology was implemented as a web-based travel assistant with
separate backend, frontend, database, local retrieval models, and external LLM
services. The separation allows the user interface and experimental runner to
reuse the same core retrieval and answer-generation services.

The backend is implemented in Python using FastAPI. SQLAlchemy provides
relational persistence, Pydantic validates API and internal pipeline schemas,
and PostgreSQL stores application data and tourism chunks. The pgvector
extension supports vector-distance operations. Sentence Transformers provides
the local embedding and CrossEncoder models. The frontend is implemented with
Angular and renders streamed Markdown responses.

\section{Backend Service}

The FastAPI application exposes grouped endpoints for authentication, chat,
memory management, health checks, and retrieval debugging. The principal API
operations include:

\begin{itemize}
    \item Google-based login, session inspection, and logout;
    \item streamed chat interaction;
    \item conversation listing and retrieval;
    \item structured memory-profile updates;
    \item memory inspection and deletion; and
    \item retrieval-debug execution for stage-level inspection.
\end{itemize}

The embedding and reranker models are loaded once during the application
lifespan and released during shutdown. This avoids loading a new model instance
for every request. The evaluation command-line runner uses the same service
functions but loads its own model registry because it executes outside the API
process.

\section{Frontend Application}

The user interface is implemented using Angular, TypeScript, PrimeNG, Tailwind
CSS, and Markdown-rendering components. It provides a conversational view in
which the user can submit a travel question and receive streamed progress and
answer content. Markdown headings, paragraphs, and lists are rendered as
structured response elements rather than displayed as raw text.

The frontend communicates with the backend through HTTP endpoints and retains
authenticated session cookies. During development, Cross-Origin Resource
Sharing is limited to the configured local Angular origins. A production
deployment should replace the development origins and cookie settings with
environment-specific secure values.

\section{Database Design}

The relational database separates application data from RAG knowledge data.
Application entities include users, conversations, messages, persistent memory
items, and session-related information. Knowledge entities include source
documents and retrieval chunks.

Each tourism chunk stores its original or cleaned text, stable hash, source
document relationship, section heading, embedding, and structured tourism
metadata. Metadata fields include geographic information, place name, place
type, chunk topic, semantic tags, activities, travel styles, suitable traveler
groups, summaries, and optional coordinates.

Persistent user memory is stored as normalized typed items. Supported types
include travel style, activity, budget, avoidance, constraint, interest,
expertise, answer length, tone, explanation style, and selected personal facts.
Each item is associated with a user and includes importance, activation state,
and timestamps. This normalized representation can be reconstructed into a
structured \texttt{UserTravelMemory} object for the retrieval and generation
pipeline.

\section{Knowledge-Base Construction Modules}

Knowledge construction is implemented as a sequence of reusable modules:

\begin{enumerate}
    \item source collection and document registration;
    \item HTML or PDF conversion to Markdown;
    \item content cleaning and duplicate removal;
    \item structure-aware Markdown chunking;
    \item LLM-assisted tourism metadata extraction;
    \item controlled-vocabulary normalization;
    \item optional geographic enrichment;
    \item embedding generation; and
    \item PostgreSQL and pgvector persistence.
\end{enumerate}

Every source and chunk retains identifiers that connect retrieved evidence to
its origin. Embedding metadata is also retained so that the corpus can be
checked for incompatible or missing vectors.

\section{Online Pipeline Services}

The online pipeline is divided into services for request routing, query
processing, memory, retrieval planning, dense retrieval, BM25 retrieval,
fusion, reranking, coverage checking, recovery, conversation-state updating,
and answer generation.

The pipeline exchanges validated Pydantic objects rather than unstructured
dictionaries wherever practical. For example, parsed constraints are normalized
to key--value objects, while locations, user memory, plans, coverage assessments,
and readiness decisions have explicit schemas. LLM responses are parsed from
JSON and validated before being accepted. Deterministic normalization repairs
known structural variations but does not silently accept arbitrary malformed
outputs.

Region normalization and destination matching combine prompt instructions with
deterministic checks. Broad areas such as Northern Vietnam are retained as
regions rather than converted to provinces. A specific city can override an
incompatible broad province value. Optional restrictive filters are relaxed
when they produce no candidates, and the relaxation is recorded.

\section{Streaming and Progress Events}

The chat endpoint streams pipeline progress so the user does not wait for the
complete answer without feedback. Events can describe query understanding,
retrieval, checking, recovery, and answer generation. Internal numerical scores
and implementation details are retained in debug traces rather than inserted
into the natural-language answer.

The final payload contains the Markdown answer, source information, route
category, and relevant identifiers. Conversation messages are saved so later
turns can resolve references. Persistent-memory extraction and temporary-state
updating occur through separate services.

\section{Retrieval Debugging Interface}

The retrieval-debug endpoint exposes the intermediate products required for
development and evaluation. These include the rewritten query, parsed schema,
planner output, applied filters, vector results, BM25 results, fused candidates,
reranker scores, selected evidence, coverage before and after recovery, filter
relaxations, missing-requirement diagnostics, answer readiness, and timing.

This interface is important because the same incorrect final response may have
different causes. For example, the parser may omit a city, filtering may remove
relevant chunks, reranking may favor one aspect, or the corpus may lack the
requested information. The trace makes these causes distinguishable.

\section{LLM Integration and Robustness}

LLM calls are wrapped by a shared telemetry layer. The wrapper records stage,
model, latency, token counts, and estimated cost when usage data is available.
Structured calls use bounded output limits and explicit JSON schemas.

If query parsing fails, a guarded fallback preserves a usable travel-information
representation rather than terminating the full request. If conversation-state
generation returns truncated or malformed JSON, the previous valid state is
kept. Coverage and recovery are bounded, preventing an unlimited loop of LLM
and retrieval calls.

The experimental pipeline is resumable. A completed case is written to the
pipeline trace before the next case starts. Later executions validate the run
and query identifiers before reusing completed cases. Judge results are also
saved per case and reused only when the run, pipeline version, model, and judge
schema are compatible.

\section{Security, Privacy, and Operational Considerations}

API keys and database credentials are provided through environment variables
and are not included in exported evaluation artifacts. Authentication sessions
use HTTP-only cookies. Production deployment should additionally enforce secure
cookies, HTTPS, restrictive CORS origins, secret rotation, rate limiting, and
access control for debugging endpoints.

Persistent memory can contain user preferences and selected personal facts.
Only non-sensitive information that materially supports future travel
assistance should be stored. The API allows memories to be inspected and
deleted. Future production use would require an explicit retention policy,
consent controls, and data-protection review.

\section{Implementation Summary}

The implementation connects the proposed multi-stage methodology to a working
application and a reproducible evaluation runner. FastAPI and Angular provide
the service and interaction layers; PostgreSQL and pgvector provide structured
and vector storage; local sentence models provide retrieval and reranking; and
external LLMs provide structured reasoning and answer generation. Shared
schemas and trace formats ensure that the deployed pipeline and experimental
pipeline can be examined using the same conceptual stages.

